\beginsong{La strasbourgeoise}
\textnote{Capo 1}
\beginverse
Petit \[Am]papa, voici la mi-carême,
\[Dm]Car te voici déguisé en sol\[Am]dat.
\[Dm]Petit papa, dis-moi si c'est pour \[Am]rire,
Ou pour \[E]faire peur aux tous petits en\[Am]fants. (bis)
\endverse
\beginverse
Non mon en\[Am]fant, je pars pour la Patrie,
C'est un de\[Dm]voir où tous les papas s'en \[Am]vont.
Embrasse-\[Dm]moi, petite fille ché\[Am]rie,
Je rentrer\[E]ai bien vite à la mai\[Am]son. (bis)
\endverse
\beginverse
Dis moi ma\[Am]man, quelle est cette médaille,
Et cette \[Dm]lettre qu'apporte le fac\[Am]teur ?
Dis moi ma\[Dm]man, tu pleures et tu dé\[Am]failles,
Ils ont tu\[E]é petit père ado\[Am]ré. (bis)
\endverse
\beginverse
Oui mon en\[Am]fant, ils ont tué ton père,
Pleurons en\[Dm]semble car nous les haï\[Am]ssons.
Quelle guerre a\[Dm]troce, qui fait pleurer les \[Am]mères,
Et tue les \[E]pères des petits anges \[Am]blonds. (bis)
\endverse
\beginverse
La neige tom\[Am]be aux portes de la ville,
Là est ass\[Dm]ise une enfant de Stras\[Am]bourg.
Elle reste \[Dm]là malgré le froid, la \[Am]bise,
Elle reste \[E]là malgré le froid du \[Am]jour. (bis)
\endverse
\beginverse
Un homme \[Am]passe, à la fillette donne,
Elle reco\[Dm]nnaît l'uniforme alle\[Am]mand.
Elle re\[Dm]fuse l'aumône qu'on lui \[Am]donne,
À l'enne\[E]mi elle dit bien fière\[Am]ment : (bis)
\endverse
\beginverse
Gardez votre \[Am]or, je garde ma puissance,
Soldat prus\[Dm]sien, passez votre che\[Am]min.
Moi je ne \[Dm]suis qu'une enfant de la \[Am]France,
À l'enne\[E]mi je ne tends pas la \[Am]main. (bis)
\endverse
\beginverse
Tout en pri\[Am]ant sous cette cathédrale,
Ma mère est \[Dm]morte sous ce porche écrou\[Am]lé.
Frappée à \[Dm]mort par l'une de vos \[Am]balles,
Frappée à \[E]mort par l'un de vos bou\[Am]lets. (bis)
\endverse
\beginverse
Mon père est \[Am]mort sur vos champs de batailles,
Je n'ai pas \[Dm]vu l'ombre de son cer\[Am]cueil.
Frappé à \[Dm]mort par l'une de vos \[Am]balles,
C'est la rai\[E]son de ma robe de \[Am]deuil. (bis)
\endverse
\beginverse
Vous avez \[Am]eu l'Alsace et la Lorraine,
Vous avez \[Dm]eu des millions d'étran\[Am]gers.
Vous avez \[Dm]eu Germanie et Bo\[Am]hème,
Mais mon p'\[E]tit c{\oe}ur il restera fran\[Am]çais. (bis)
\endverse
\endsong
